\documentclass[aspectratio=169,xcolor=dvipsnames]{beamer}

\AtBeginSection[]
{
    \begin{frame}
        \frametitle{Visão Geral}
        \tableofcontents[currentsection]
    \end{frame}
}

\usefonttheme{professionalfonts} % using non standard fonts for beamer
\usefonttheme{serif} % default family is serif
\usepackage{fontspec}
\setmainfont{XCharter}[
    UprightFont    = *-Roman,
    BoldFont       = *-Bold,
    ItalicFont     = *-Italic,
    BoldItalicFont = *-BoldItalic,
]

\usepackage{hyperref}
\usepackage{graphicx} % Allows including images
\usepackage{booktabs} % Allows the use of \toprule, \midrule and \bottomrule in tables
\input{./slides/SimplePlusTheme/beamerthemeSimplePlus.sty}
\input{./slides/SimplePlusTheme/beamerinnerthemeSimplePlus.sty}
\input{./slides/SimplePlusTheme/beamerfontthemeSimplePlus.sty}
\input{./slides/SimplePlusTheme/beamercolorthemeSimplePlus.sty}

\usepackage{listings}
\setbeamercovered{transparent}

\usepackage{biblatex}
\addbibresource{slides/reference.bib}
\renewcommand*{\bibfont}{\footnotesize}
% Tamanho dos slides de "Leituras Recomendadas": maior que a bibliografia
% completa (\bibfont acima, que rende \tiny via a redefinição de \footnotesize)
% para destacar as leituras curadas.
\newcommand{\leiturasize}{\small}

\usepackage[brazilian ]{babel}
\usepackage{csquotes}
\usepackage{pgfplots}
\pgfplotsset{compat=1.18}

\usepackage{emoji}
\usepackage{amsmath}
\usepackage[xcharter,bigdelims,vvarbb]{newtxmath}
% newtxmath's operators font requests the fontspec family `XCharter(0)' in OT1;
% without these substitutions math digits and operator names silently fall back
% to Computer Modern (LaTeX Font Warning: Font shape `OT1/XCharter(0)/m/n' undefined).
\DeclareFontFamily{OT1}{XCharter(0)}{}
\DeclareFontShape{OT1}{XCharter(0)}{m}{n}{<->ssub * XCharter-TLF/m/n}{}
\DeclareFontShape{OT1}{XCharter(0)}{m}{it}{<->ssub * XCharter-TLF/m/it}{}
\DeclareFontShape{OT1}{XCharter(0)}{b}{n}{<->ssub * XCharter-TLF/b/n}{}
\DeclareFontShape{OT1}{XCharter(0)}{bx}{n}{<->ssub * XCharter-TLF/b/n}{}

\usepackage{bm}
\usepackage[table]{xcolor}
\usepackage{xcolor}
\usepackage{algpseudocode}

\usepackage{cancel}
\usepackage{ulem}

\usetikzlibrary{shapes}
\usetikzlibrary{arrows}
\usetikzlibrary {arrows.meta}
\usetikzlibrary{calc,positioning}
\usepgfplotslibrary{fillbetween}
\usetikzlibrary{angles,quotes}
\usetikzlibrary{tikzmark}


\definecolor{PastelRed}{HTML}{FFC1CC}
\definecolor{PastelBlue}{HTML}{C2F0FF}
\definecolor{PastelBlue2}{HTML}{3182BD}

\definecolor{PastelPink}{HTML}{FFCBE1}
\definecolor{PastelGreen}{HTML}{D6E5BD}
\definecolor{PastelYellow}{HTML}{F9E1A8}
\definecolor{PastelBlue}{HTML}{BCD8EC}
\definecolor{PastelPurple}{HTML}{DCCCEC}
\definecolor{PastelOrange}{HTML}{FFDAB4}


\renewcommand{\footnotesize}{\tiny}

\newcommand{\mespor}{
  \ifcase\month
    \or janeiro
    \or fevereiro
    \or março
    \or abril
    \or maio
    \or junho
    \or julho
    \or agosto
    \or setembro
    \or outubro
    \or novembro
    \or dezembro
  \fi
}
% \usepackage{csquotes}
% \usepackage{minted}
% \usemintedstyle{xcode}
\definecolor{vscLightBackground}{HTML}{FFFFFF}
\definecolor{vscLightText}{HTML}{000000}
\definecolor{vscLightGreen}{HTML}{008000}        % For Comments AND Numbers
\definecolor{vscLightRed}{HTML}{A31515}          % For Strings
\definecolor{vscLightTeal}{HTML}{267F99}
\definecolor{vscBlue}{HTML}{0000FF}% For main keywords (if, for, def...)
\definecolor{vscLightPytorch}{HTML}{800080}      % For PyTorch keywords (dark magenta)

\lstdefinestyle{vscode-light}{
    language=Python,
    backgroundcolor=\color{vscLightBackground},
    basicstyle=\ttfamily\scriptsize\color{vscLightText},
    keywordstyle=\color{vscLightTeal},
    commentstyle=\color{vscLightGreen},
    stringstyle=\color{vscLightRed},
    % Style for PyTorch keywords (using keyword group 2)
    keywordstyle=[2]\color{vscLightPytorch}, 
    morekeywords={self, True, False, None},
    morekeywords=[2]{
        torch, Tensor, device, dtype, to, nn, optim, autograd, utils, data,
        Module, Linear, Conv2d, ReLU, CrossEntropyLoss, Adam, SGD,
        Dataset, DataLoader, randn, zeros, ones, arange, cat, stack,
        view, reshape, sum, mean, max, min, matmul, no_grad, backward
    },
    keywordstyle=[3]\color{vscBlue}, 
    morekeywords=[3]{
        model, loss, y, y_hat, x , optimizer
    },
    frame=single,
    tabsize=1,
    showstringspaces=false,
    breaklines=true,
    captionpos=b,
    extendedchars=true,
    numbers=left,
}

\usepackage[cache=true]{minted}
% minted v3+ (TeX Live 2024+) auto-detects the output directory.
% minted v2 (Ubuntu apt TeX Live 2023) needs it set explicitly to
% match latexmkrc's $out_dir = 'build'.
\makeatletter
\@ifpackagelater{minted}{2024/01/01}{}{%
  \def\minted@outputdir{build/}%
}
\makeatother
\setminted{
    fontsize=\scriptsize,
    style=vs,
    breaklines=true,
    numbers=left,
    numbersep=5pt,
    frame=single,
}
\providecommand{\citeauthoryear}[1]{(\citeauthor{#1},~\citeyear{#1})}

\providecommand{\thetav}{\bm{\theta}}

% vetor coluna 2D com colchetes
\providecommand{\cvec}[2]{\begin{bmatrix} #1 \\ #2 \end{bmatrix}}

\providecommand{\varvector}[1]{\mathbf{#1}}
% instance
\providecommand{\X}{\varvector{X}}
\providecommand{\mlinstance}{\varvector{X}}
\providecommand{\mlinstancei}{\varvector{X}^{(i)}}
\providecommand{\mlinstancex}[1]{\varvector{X}^{(#1)}}

% target
\providecommand{\y}{\varvector{y}}
\providecommand{\mltarget}{\varvector{y}}
\providecommand{\mltargeti}{\varvector{y}^{(i)}}
\providecommand{\mltargetx}[1]{\varvector{y}^{(#1)}}

%hats
\providecommand{\yhat}{\hat{\varvector{y}}}
\providecommand{\mltargethat}{\hat{\varvector{y}}}
\providecommand{\mltargethati}{\hat{\varvector{y}}^{(i)}}
\providecommand{\mltargethatx}[1]{\hat{\varvector{y}}^{(#1)}}

% linear reg
\providecommand{\mllinearregression}{\mltargethat=\mlinstance\bm{\theta}}

\providecommand{\tikzarrowcolor}[3]{%
  \begin{tikzpicture}[remember picture, overlay]
    % Arrow body
    \path[fill=#3] ([shift={(#1,#2)}]current page.south west) 
      rectangle ++(0.3, -0.2);
    % Arrow head
    \path[fill=#3] ([shift={(#1+0.3,#2+0.1)}]current page.south west) -- 
         ++(0, -0.4) -- 
         ++(0.2, 0.2) -- 
         cycle;
  \end{tikzpicture}%
}

\providecommand{\tikzarrow}[2]{%
    \tikzarrowcolor{#1}{#2}{Red}
}



\DeclareMathOperator*{\argmax}{argmax}
\DeclareMathOperator*{\argmin}{argmin}



% Tokenization-specific colors (used by \tok commands below)
\definecolor{pucrsblue}{RGB}{0, 60, 113}
\definecolor{lightblue}{RGB}{220, 235, 252}
\definecolor{lightgreen}{RGB}{220, 245, 220}
\definecolor{lightyellow}{RGB}{255, 252, 220}
\definecolor{lightred}{RGB}{255, 225, 225}
\definecolor{lightorange}{RGB}{255, 237, 210}
\definecolor{darkgreen}{RGB}{0, 120, 0}

% Extra packages needed for this lecture
\usepackage{array}
\usepackage{multirow}

% fontawesome5 not available; using text symbols instead
\newcommand{\faCheck}{\textcolor{darkgreen}{\checkmark}}
\newcommand{\faTimes}{\textcolor{red}{$\times$}}

% --- Custom commands ---
\newcommand{\tok}[1]{\fcolorbox{pucrsblue}{lightblue}{\strut\textcolor{black}{\texttt{#1}}}}
\newcommand{\toka}[1]{\fcolorbox{darkgreen}{lightgreen}{\strut\textcolor{black}{\texttt{#1}}}}
\newcommand{\tokb}[1]{\fcolorbox{red!70!black}{lightred}{\strut\textcolor{black}{\texttt{#1}}}}
\newcommand{\tokc}[1]{\fcolorbox{orange!80!black}{lightorange}{\strut\textcolor{black}{\texttt{#1}}}}
\newcommand{\mergetok}[1]{\fcolorbox{orange!80!black}{lightyellow}{\strut\textcolor{black}{\textbf{\texttt{#1}}}}}

% --- Title ---
\title{Deep Learning}
\subtitle{Tokenização em Modelos de Linguagem}
\author{Lucas Silveira Kupssinskü}
\institute{Machine Learning Theory and Applications Laboratory \\ PUCRS}
\date{\mespor de \the\year}

\begin{document}

% ============================================================
\begin{frame}
\titlepage
\end{frame}

% ============================================================
\begin{frame}{Visão Geral}
\tableofcontents
\end{frame}

% ============================================================
\section{Por que a Tokenization Importa}
% ============================================================

\begin{frame}{O que é Tokenization?}

Tokenization é o processo de converter texto bruto em uma sequência de unidades discretas (\textbf{tokens}) que um modelo pode processar.

\vspace{6pt}

\begin{block}{O tokenizer define:}
\begin{enumerate}
    \item O \textbf{vocabulário}: quais tokens existem
    \item A \textbf{segmentação}: como o texto é dividido em tokens
    \item O \textbf{mapeamento}: tokens $\leftrightarrow$ IDs inteiros
\end{enumerate}
\end{block}

\vspace{6pt}

\centering
\begin{tikzpicture}[node distance=0.4cm, >=Stealth, font=\small]
    \node[draw, rounded corners, fill=lightyellow, minimum width=3.5cm, minimum height=0.7cm] (text) {\texttt{"Hello world"}};
    \node[draw, rounded corners, fill=lightblue, minimum width=3.5cm, minimum height=0.7cm, right=1.5cm of text] (tokens) {\tok{Hello} \tok{world}};
    \node[draw, rounded corners, fill=lightgreen, minimum width=3.5cm, minimum height=0.7cm, right=1.5cm of tokens] (ids) {\texttt{[15496, 995]}};
    \draw[->, thick] (text) -- node[above]{\scriptsize tokenize} (tokens);
    \draw[->, thick] (tokens) -- node[above]{\scriptsize encode} (ids);
\end{tikzpicture}

\end{frame}

% ============================================================

\begin{frame}{Por que isso importa?}

A tokenization \textbf{molda silenciosamente} tudo que vem a seguir:

\vspace{8pt}

\begin{columns}[T]
\column{0.48\textwidth}
\begin{alertblock}{Problemas causados pela tokenization}
\begin{itemize}
    \item O GPT não consegue contar letras de forma confiável em \texttt{"strawberry"} (3 r's)
    \item Desigualdade multilíngue: o inglês usa menos tokens que outros idiomas
    \item \texttt{"123 + 456"} pode ser dividido de forma inconsistente entre os dígitos
\end{itemize}
\end{alertblock}

\column{0.48\textwidth}
\begin{exampleblock}{Trade-offs de design}
\begin{itemize}
    \item Vocabulário pequeno $\rightarrow$ sequências longas
    \item Vocabulário grande $\rightarrow$ mais parâmetros e tokens raros subtreinados
    \item Dificilmente encontraremos uma ``granularidade correta''
\end{itemize}
\end{exampleblock}
\end{columns}

\vspace{8pt}
\centering
\textbf{O tokenizer não é um detalhe de pré-processamento; é uma decisão arquitetural.}

\end{frame}

% ============================================================

\begin{frame}{O Problema da Fertility}

\textbf{Fertility} = número de tokens necessários para representar uma palavra ou frase.

\vspace{6pt}

\begin{center}
\renewcommand{\arraystretch}{1.3}
\begin{tabular}{lcc}
\toprule
\textbf{Idioma} & \textbf{Frase} & \textbf{Tokens GPT-2} \\
\midrule
Inglês   & ``Hello, how are you?'' & 6 tokens \\
Espanhol & ``Hola, \textquestiondown c\'omo est\'as?'' & 9 tokens \\
Japonês  & (equivalente) & 14+ tokens \\
\bottomrule
\end{tabular}
\end{center}

\vspace{6pt}

\begin{alertblock}{Consequência}
Maior fertility $=$ mais tokens $=$ mais custo computacional, janela de contexto efetiva menor, e frequentemente pior desempenho para esse idioma.
\end{alertblock}

\end{frame}

% ============================================================
\section{Abordagens Ingênuas}
% ============================================================

\begin{frame}{Nosso Exemplo Recorrente}

Ao longo desta aula usamos o mesmo \textbf{corpus de treinamento}:

\vspace{6pt}

\begin{center}
\renewcommand{\arraystretch}{1.3}
\begin{tabular}{lc}
\toprule
\textbf{Palavra} & \textbf{Frequência no corpus} \\
\midrule
\texttt{low}    & 5 \\
\texttt{lower}  & 2 \\
\texttt{newest} & 6 \\
\texttt{widest} & 3 \\
\bottomrule
\end{tabular}
\end{center}

\vspace{6pt}

\begin{block}{Por que as frequências importam}
Um tokenizer é \textbf{aprendido} a partir de um corpus. A frequência de cada palavra
(e de cada par de caracteres) determina quais tokens são adicionados ao vocabulário.
\end{block}

\end{frame}

% ============================================================

\begin{frame}{Tokenization em Nível de Caractere}

\textbf{Ideia:} Cada caractere é um token.

\vspace{4pt}

\begin{exampleblock}{Exemplo Recorrente}
\texttt{"low lower newest widest"} \\[4pt]
$\rightarrow$ \tok{l}\tok{o}\tok{w}\tok{\_}~\tok{l}\tok{o}\tok{w}\tok{e}\tok{r}\tok{\_}~\tok{n}\tok{e}\tok{w}\tok{e}\tok{s}\tok{t}\tok{\_}~\tok{w}\tok{i}\tok{d}\tok{e}\tok{s}\tok{t} \\[4pt]
Vocabulário: \{~\texttt{l, o, w, \_, e, r, n, s, t, i, d}~\} $\quad|\mathcal{V}| = 11$
\end{exampleblock}

\vspace{4pt}

\begin{columns}[T]
\column{0.48\textwidth}
\textcolor{darkgreen}{\faCheck~\textbf{Vantagens}}
\begin{itemize}
    \item Vocabulário reduzido ($\sim$256 para ASCII)
    \item Sem palavras fora do vocabulário (OOV)
    \item Agnóstico ao idioma
\end{itemize}

\column{0.48\textwidth}
\textcolor{red}{\faTimes~\textbf{Desvantagens}}
\begin{itemize}
    \item Sequências muito longas
    \item Perde a semântica em nível de palavra
    \item O modelo precisa ``aprender'' a ortografia
\end{itemize}
\end{columns}

\end{frame}

% ============================================================

\begin{frame}{Tokenization em Nível de Palavra}

\textbf{Ideia:} Dividir por espaço em branco (e pontuação). Cada palavra é um token.

\vspace{2pt}

\begin{exampleblock}{Exemplo Recorrente}
\texttt{"low lower newest widest"} \\[4pt]
$\rightarrow$ \tok{low} \tok{lower} \tok{newest} \tok{widest} \\[4pt]
Vocabulário: \{~\texttt{low, lower, newest, widest}~\} $\quad|\mathcal{V}| = 4$
\end{exampleblock}

\vspace{4pt}

\begin{columns}[T]
\column{0.48\textwidth}
\textcolor{darkgreen}{\faCheck~\textbf{Vantagens}}
\begin{itemize}
    \item Sequências curtas
    \item Tokens carregam significado semântico
    \item Simples de implementar
\end{itemize}

\column{0.48\textwidth}
\textcolor{red}{\faTimes~\textbf{Desvantagens}}
\begin{itemize}
    \item Vocabulário gigantesco (100K+)
    \item Não lida com palavras OOV
    \item Sem compartilhamento morfológico: \texttt{low} e \texttt{lowest} são tokens sem relação
\end{itemize}
\end{columns}

\vspace{2pt}
\centering
\textbf{A sacada principal:} Queremos algo \textit{entre} caracteres e palavras.\\
$\Rightarrow$ \textbf{Tokenization por subwords}

\end{frame}

% ============================================================
\section{Byte Pair Encoding (BPE)}
% ============================================================

\begin{frame}{BPE: Ideia Central}

\textbf{Byte Pair Encoding} (Sennrich et al., 2016) mescla iterativamente o par de símbolos adjacentes mais frequente.

\vspace{6pt}

\begin{block}{Algoritmo}
\begin{enumerate}
    \item Começa com um vocabulário de todos os caracteres individuais
    \item Conta todos os pares de símbolos adjacentes no corpus
    \item Mescla o par mais frequente em um novo símbolo
    \item Repete os passos 2--3 por $k$ merges (hiperparâmetro)
\end{enumerate}
\end{block}

\vspace{6pt}

\textbf{Propriedades principais:}
\begin{itemize}
    \item Bottom-up, guloso, baseado em frequência
    \item A tabela de merges é aprendida a partir de um corpus de treinamento
    \item Codificação determinística na inferência (aplica os merges na ordem aprendida)
\end{itemize}

\vspace{4pt}
\centering
\textbf{Utilizado em:} GPT, GPT-2, GPT-3, RoBERTa, CodeGen, \ldots

\end{frame}

% ============================================================

\begin{frame}{BPE: Configuração do Exemplo Simplificado}

\textbf{Corpus de treinamento} (com frequências de palavras e marcador de fim de palavra \texttt{\_}):

\vspace{4pt}

\begin{center}
\renewcommand{\arraystretch}{1.2}
\begin{tabular}{lcl}
\toprule
\textbf{Palavra} & \textbf{Freq} & \textbf{Divisão inicial} \\
\midrule
\texttt{low\_}    & 5 & \tok{l}~\tok{o}~\tok{w}~\tok{\_} \\
\texttt{lower\_}  & 2 & \tok{l}~\tok{o}~\tok{w}~\tok{e}~\tok{r}~\tok{\_} \\
\texttt{newest\_} & 6 & \tok{n}~\tok{e}~\tok{w}~\tok{e}~\tok{s}~\tok{t}~\tok{\_} \\
\texttt{widest\_} & 3 & \tok{w}~\tok{i}~\tok{d}~\tok{e}~\tok{s}~\tok{t}~\tok{\_} \\
\bottomrule
\end{tabular}
\end{center}

\vspace{6pt}

\textbf{Vocabulário inicial:}
$\mathcal{V} = \{$ \texttt{l, o, w, \_, e, r, n, s, t, i, d} $\}$ \quad $|\mathcal{V}| = 11$

\vspace{6pt}
Vamos executar os merges passo a passo\ldots

\end{frame}

% ============================================================

\begin{frame}[shrink=20]{BPE: Merge Passo 1}

\textbf{Contagem de todos os pares adjacentes:}

\vspace{2pt}
\begin{center}
\small
\renewcommand{\arraystretch}{1.1}
\begin{tabular}{lclclclc}
\toprule
\textbf{Par} & \textbf{Contagem} & \textbf{Par} & \textbf{Contagem} & \textbf{Par} & \textbf{Contagem} & \textbf{Par} & \textbf{Contagem} \\
\midrule
(e, s) & \textbf{9} & (l, o) & 7 & (o, w) & 7 & (s, t) & 9 \\
(e, w) & 6 & (w, \_) & 5 & (t, \_) & 9 & (e, r) & 2 \\
(w, e) & 8 & (n, e) & 6 & (w, i) & 3 & (d, e) & 3 \\
(i, d) & 3 & (r, \_) & 2 & & & & \\
\bottomrule
\end{tabular}
\end{center}

\vspace{4pt}

Mais frequentes: \textbf{(e, s)}, \textbf{(s, t)} e \textbf{(t, \_)} têm contagem 9. \\
\textit{(Desempate: escolher (e, s) primeiro.)}

\vspace{4pt}

\begin{exampleblock}{Merge 1: (e, s) $\rightarrow$ \mergetok{es} \quad (contagem: 9)}
\begin{tabular}{lcl}
\texttt{low\_}    & 5 & \tok{l}~\tok{o}~\tok{w}~\tok{\_} \\
\texttt{lower\_}  & 2 & \tok{l}~\tok{o}~\tok{w}~\tok{e}~\tok{r}~\tok{\_} \\
\texttt{newest\_} & 6 & \tok{n}~\tok{e}~\tok{w}~\mergetok{es}~\tok{t}~\tok{\_} \\
\texttt{widest\_} & 3 & \tok{w}~\tok{i}~\tok{d}~\mergetok{es}~\tok{t}~\tok{\_} \\
\end{tabular}
\end{exampleblock}

$\mathcal{V} = \{$ \texttt{l, o, w, \_, e, r, n, s, t, i, d, \textbf{es}} $\}$ \quad $|\mathcal{V}| = 12$

\end{frame}

% ============================================================

\begin{frame}[shrink=15]{BPE: Merges Passos 2--4}

\begin{exampleblock}{Merge 2: (es, t) $\rightarrow$ \mergetok{est} \quad (contagem: 9)}
\begin{tabular}{lcl}
\texttt{newest\_} & 6 & \tok{n}~\tok{e}~\tok{w}~\mergetok{est}~\tok{\_} \\
\texttt{widest\_} & 3 & \tok{w}~\tok{i}~\tok{d}~\mergetok{est}~\tok{\_} \\
\end{tabular}
\quad\quad (demais inalterados)
\end{exampleblock}

\begin{exampleblock}{Merge 3: (est, \_) $\rightarrow$ \mergetok{est\_} \quad (contagem: 9)}
\begin{tabular}{lcl}
\texttt{newest\_} & 6 & \tok{n}~\tok{e}~\tok{w}~\mergetok{est\_} \\
\texttt{widest\_} & 3 & \tok{w}~\tok{i}~\tok{d}~\mergetok{est\_} \\
\end{tabular}
\end{exampleblock}

\begin{exampleblock}{Merge 4: (l, o) $\rightarrow$ \mergetok{lo} \quad (contagem: 7)}
\begin{tabular}{lcl}
\texttt{low\_}    & 5 & \mergetok{lo}~\tok{w}~\tok{\_} \\
\texttt{lower\_}  & 2 & \mergetok{lo}~\tok{w}~\tok{e}~\tok{r}~\tok{\_} \\
\end{tabular}
\end{exampleblock}

$\mathcal{V} = \{$ \texttt{l, o, w, \_, e, r, n, s, t, i, d, es, est, est\_, lo} $\}$ \quad $|\mathcal{V}| = 15$

\end{frame}

% ============================================================

\begin{frame}[shrink=10]{BPE: Merges Passos 5--7}

\begin{exampleblock}{Merge 5: (lo, w) $\rightarrow$ \mergetok{low} \quad (contagem: 7)}
\begin{tabular}{lcl}
\texttt{low\_}    & 5 & \mergetok{low}~\tok{\_} \\
\texttt{lower\_}  & 2 & \mergetok{low}~\tok{e}~\tok{r}~\tok{\_} \\
\end{tabular}
\end{exampleblock}

\begin{exampleblock}{Merge 6: (n, e) $\rightarrow$ \mergetok{ne} \quad (contagem: 6)}
\begin{tabular}{lcl}
\texttt{newest\_} & 6 & \mergetok{ne}~\tok{w}~\tok{est\_} \\
\end{tabular}
\end{exampleblock}

\begin{exampleblock}{Merge 7: (ne, w) $\rightarrow$ \mergetok{new} \quad (contagem: 6)}
\begin{tabular}{lcl}
\texttt{newest\_} & 6 & \mergetok{new}~\tok{est\_} \\
\end{tabular}
\end{exampleblock}

\vspace{4pt}

Após 7 merges, o corpus é representado como:

\begin{center}
\tok{low}~\tok{\_} ~~(5) \quad\quad
\tok{low}~\tok{e}~\tok{r}~\tok{\_} ~~(2) \quad\quad
\tok{new}~\tok{est\_} ~~(6) \quad\quad
\tok{w}~\tok{i}~\tok{d}~\tok{est\_} ~~(3)
\end{center}

Observe como \tok{est\_} é compartilhado entre ``newest'' e ``widest''!

\end{frame}

% ============================================================

\begin{frame}[shrink=25]{BPE: Codificando Novo Texto}

Na inferência, aplica-se os merges aprendidos \textbf{em ordem}:

\vspace{6pt}

\begin{exampleblock}{Codificando a palavra ``lowest''}
\begin{enumerate}
    \item Início: \tok{l}~\tok{o}~\tok{w}~\tok{e}~\tok{s}~\tok{t}~\tok{\_}
    \item Aplicar merge 1 (e,s)$\rightarrow$es: \tok{l}~\tok{o}~\tok{w}~\tok{es}~\tok{t}~\tok{\_}
    \item Aplicar merge 2 (es,t)$\rightarrow$est: \tok{l}~\tok{o}~\tok{w}~\tok{est}~\tok{\_}
    \item Aplicar merge 3 (est,\_)$\rightarrow$est\_: \tok{l}~\tok{o}~\tok{w}~\tok{est\_}
    \item Aplicar merge 4 (l,o)$\rightarrow$lo: \tok{lo}~\tok{w}~\tok{est\_}
    \item Aplicar merge 5 (lo,w)$\rightarrow$low: \tok{low}~\tok{est\_}
\end{enumerate}
\end{exampleblock}

\vspace{4pt}

\textbf{Resultado:} \tok{low}~\tok{est\_} compartilha subwords com ``low'', ``newest'', ``widest''!

\vspace{4pt}

Mesmo que ``lowest'' \textbf{nunca tenha aparecido no corpus de treinamento}, o modelo consegue compô-la a partir de subwords aprendidos.

\end{frame}

% ============================================================

\begin{frame}[fragile,shrink=25]{BPE: Esboço em Python}
\begin{lstlisting}[style=vscode-light, numbers=none]
import re, collections

def get_stats(vocab):
    """Conta a frequência de pares de símbolos adjacentes."""
    pairs = collections.Counter()
    for word, freq in vocab.items():
        symbols = word.split()
        for i in range(len(symbols) - 1):
            pairs[(symbols[i], symbols[i+1])] += freq
    return pairs

def merge_vocab(pair, vocab):
    """Mescla todas as ocorrências de um par no vocabulário."""
    bigram = re.escape(' '.join(pair))
    pattern = re.compile(r'(?<!\S)' + bigram + r'(?!\S)')
    new_vocab = {}
    for word in vocab:
        new_word = pattern.sub(''.join(pair), word)
        new_vocab[new_word] = vocab[word]
    return new_vocab

# Corpus de treinamento (frequências de palavras)
vocab = {'l o w _': 5, 'l o w e r _': 2,
         'n e w e s t _': 6, 'w i d e s t _': 3}

num_merges = 7
for i in range(num_merges):
    pairs = get_stats(vocab)
    best = max(pairs, key=pairs.get)
    vocab = merge_vocab(best, vocab)
    print(f"Merge {i+1}: {best} -> {''.join(best)}")
\end{lstlisting}
\end{frame}

% ============================================================
\section{WordPiece}
% ============================================================

\begin{frame}{WordPiece: Ideia Central}

\textbf{WordPiece} (Schuster \& Nakajima, 2012; Wu et al., 2016) é similar ao BPE, mas seleciona merges com base na \textbf{likelihood} em vez da frequência.

\vspace{3pt}

\begin{block}{Diferença principal em relação ao BPE}
Em vez de mesclar o par mais \textit{frequente}, o WordPiece mescla o par que \textbf{maximiza a likelihood dos dados de treinamento}:
$$
\text{score}(x, y) = \frac{\text{freq}(xy)}{\text{freq}(x) \cdot \text{freq}(y)}
$$
Isso favorece mesclar pares cuja combinação é muito mais provável do que o acaso.
\end{block}

\vspace{2pt}

\textbf{Outras diferenças:}
\begin{itemize}
    \item Usa o prefixo \texttt{\#\#} para tokens de continuação: ``playing'' $\rightarrow$ \tok{play}~\tok{\#\#ing}
    \item Longest-match guloso da esquerda para a direita na inferência
\end{itemize}

\vspace{2pt}
\centering
\textbf{Utilizado em:} BERT, DistilBERT, Electra, \ldots

\end{frame}

% ============================================================

\begin{frame}{WordPiece: Exemplo Simplificado}

\textbf{Mesmo corpus, scores diferentes:}

\vspace{4pt}

\begin{center}
\small
\renewcommand{\arraystretch}{1.2}
\begin{tabular}{lcccl}
\toprule
\textbf{Par} & \textbf{freq(xy)} & \textbf{freq(x)$\cdot$freq(y)} & \textbf{Score} & \\
\midrule
(e, s) & 9 & $17 \times 9 = 153$  & 0.0588 & \\
(s, t) & 9 & $9 \times 9 = 81$   & 0.111  & \\
(i, d) & 3 & $3 \times 3 = 9$    & \textbf{0.333} & $\leftarrow$ highest \\
(n, e) & 6 & $6 \times 17 = 102$ & 0.0588 & \\
(l, o) & 7 & $7 \times 7 = 49$   & 0.143  & \\
\bottomrule
\end{tabular}
\end{center}

\vspace{4pt}

\textbf{O WordPiece mescla (i, d) primeiro}, não (e, s)!

\vspace{2pt}

\begin{alertblock}{Intuição}
(i, d) aparecem quase \textit{exclusivamente} juntos $\rightarrow$ alta informação mútua. \\
(e, s) aparecem com frequência, mas também de forma independente $\rightarrow$ score menor.
\end{alertblock}

\end{frame}

% ============================================================

\begin{frame}{WordPiece: Inferência: Longest Match}

Na inferência, o WordPiece usa \textbf{longest-match-first guloso} da esquerda para a direita.

\vspace{6pt}

\begin{exampleblock}{Exemplo: codificando ``unhappiness'' com um vocabulário estilo BERT}

Suponha que o vocabulário contém: \texttt{un}, \texttt{happy}, \texttt{\#\#happi}, \texttt{\#\#ness}, \texttt{\#\#ly}, \ldots

\vspace{4pt}

\begin{enumerate}
    \item Varredura da esquerda: maior correspondência = \tok{un}
    \item Restante: ``happiness''. Maior correspondência = \tok{\#\#happi}
    \item Restante: ``ness''. Maior correspondência = \tok{\#\#ness}
\end{enumerate}

\vspace{4pt}
Resultado: \tok{un}~\tok{\#\#happi}~\tok{\#\#ness}
\end{exampleblock}

\vspace{6pt}

O prefixo \texttt{\#\#} indica ``este token é uma continuação da palavra anterior'' (não é o início de uma palavra).

\end{frame}

% ============================================================
\section{Unigram Language Model (SentencePiece)}
% ============================================================

\begin{frame}{Unigram: Ideia Central}

\textbf{Unigram LM} (Kudo, 2018) adota a \textbf{abordagem oposta} à do BPE:

\vspace{6pt}

\begin{block}{Algoritmo}
\begin{enumerate}
    \item Começa com um vocabulário inicial \textbf{grande} (ex.: todos os substrings até certo comprimento)
    \item Define um modelo de linguagem unigrama: $P(x) = \prod_{i=1}^{n} P(x_i)$
    \item Para cada token, calcula o quanto a likelihood \textbf{cai} se ele for removido
    \item Remove os tokens com a \textbf{menor perda} (menos úteis)
    \item Repete até o vocabulário atingir o tamanho desejado
\end{enumerate}
\end{block}

\vspace{6pt}

\begin{alertblock}{Diferença principal em relação ao BPE}
BPE: começa pequeno, cresce (bottom-up) \\
Unigram: começa grande, poda (top-down)
\end{alertblock}

\end{frame}

% ============================================================

\begin{frame}{Unigram: Exemplo Simplificado (1/2)}

\textbf{Corpus:} ``low'' (freq 10), ``er'' (freq 10), ``lower'' (freq 1)

\vspace{6pt}

\textbf{Passo 1:} Começa com vocabulário de todos os substrings; $P(t) = \tfrac{\text{count}(t)}{N}$, onde $N = \sum_{t'}\text{count}(t')$:

\begin{center}
\small
\renewcommand{\arraystretch}{1.1}
\begin{tabular}{lclclc}
\toprule
\textbf{Token} & $P(\cdot)$ & \textbf{Token} & $P(\cdot)$ & \textbf{Token} & $P(\cdot)$ \\
\midrule
\texttt{low}   & 0.105 & \texttt{lo}   & 0.105 & \texttt{l} & 0.105 \\
\texttt{er}    & 0.105 & \texttt{ow}   & 0.105 & \texttt{o} & 0.105 \\
\texttt{lower} & 0.010 & \texttt{we}   & 0.010 & \texttt{w} & 0.105 \\
\texttt{owe}   & 0.010 & \texttt{wer}  & 0.010 & \texttt{e} & 0.105 \\
\texttt{lowe}  & 0.010 & \texttt{ower} & 0.010 & \texttt{r} & 0.105 \\
\bottomrule
\end{tabular}
\end{center}

{\footnotesize Substrings de ``low'' ou de ``er'' têm count $= 11$ ($= 10$ ocorrências da sua palavra $+ 1$ de ``lower''); substrings exclusivas de ``lower'' têm count $= 1$; $N = 9{\times}11 + 6{\times}1 = 105$.}

\end{frame}

% ============================================================

\begin{frame}{Unigram: Exemplo Simplificado (2/2)}

\textbf{Corpus:} ``low'' (freq 10), ``er'' (freq 10), ``lower'' (freq 1)

\vspace{6pt}

\textbf{Passo 2:} Melhor segmentação de ``lower'' sob este modelo:

\vspace{4pt}
\hspace{1cm}$P(\texttt{low}) \times P(\texttt{er}) = 0.105 \times 0.105 \approx \mathbf{0.011}$ \quad $\leftarrow$ melhor \\
\hspace{1cm}$P(\texttt{lower}) = 0.010$ \quad (token único; ligeiramente inferior ao par) \\
\hspace{1cm}$P(\texttt{l}) \times P(\texttt{owe}) \times P(\texttt{r}) = 0.105 \times 0.010 \times 0.105 \approx 0.0001$

\vspace{10pt}

\textbf{Passo 3:} Poda do vocabulário:

\vspace{4pt}
\hspace{1cm}$P(\texttt{owe}) = 0.010$, nunca entra na melhor segmentação $\rightarrow$ podá-lo. \\[4pt]
\hspace{1cm}$P(\texttt{low}) = 0.105$, essencial em 11 segmentações $\rightarrow$ mantê-lo.

\end{frame}


% ============================================================

\begin{frame}{SentencePiece: O Framework}

\textbf{SentencePiece} (Kudo \& Richardson, 2018) é um \textbf{framework}, não um algoritmo.

\vspace{6pt}

\begin{block}{O que torna o SentencePiece especial?}
\begin{enumerate}
    \item \textbf{Agnóstico ao idioma:} trata a entrada como um fluxo bruto de bytes/unicode
    \item \textbf{Sem pre-tokenization:} não depende de divisão por espaço em branco
    \item O espaço em branco é tratado como um caractere regular (substituído por \texttt{\textunderscore})
    \item Suporta \textbf{tanto} BPE quanto Unigram como algoritmo subjacente
\end{enumerate}
\end{block}

\vspace{6pt}

\begin{exampleblock}{Exemplo}
\texttt{"New York"} $\rightarrow$ \tok{\_New}~\tok{\_York}

O \texttt{\_} codifica o espaço, tornando a tokenization \textbf{totalmente reversível}.
\end{exampleblock}

\vspace{4pt}
\centering
\textbf{Utilizado em:} T5, LLaMA, ALBERT, XLNet, Gemma, \ldots

\end{frame}

% ============================================================
\section{Byte-Level BPE}
% ============================================================

\begin{frame}[shrink=10]{Byte-Level BPE: Motivação}

\textbf{Problema:} O BPE padrão opera sobre caracteres Unicode. Mas o Unicode tem 150K+ pontos de código; o vocabulário base é enorme, e scripts raros geram muitos tokens \texttt{<unk>}.

\vspace{6pt}

\begin{block}{Solução: Byte-Level BPE (Radford et al., 2019, GPT-2)}
\begin{itemize}
    \item Opera sobre \textbf{bytes brutos} (sempre 256 tokens base)
    \item Todo texto pode ser codificado: \textbf{nunca há tokens desconhecidos}
    \item Em seguida, aplica-se os merges padrão do BPE sobre os bytes
\end{itemize}
\end{block}

\vspace{6pt}

\begin{exampleblock}{Exemplo}
O emoji caveira ``\emoji{skull}'' (\texttt{U+1F480}) ocupa 4 bytes em UTF-8, \texttt{0xF0 0x9F 0x92 0x80}:

\vspace{2pt}
\hspace{1cm} $\rightarrow$ \tok{0xF0}~\tok{0x9F}~\tok{0x92}~\tok{0x80} \quad (4 tokens de byte)

\vspace{2pt}
Mesmo fora do vocabulário, ele \textbf{nunca vira \texttt{<unk>}}: todo caractere é representável byte a byte, e bytes frequentes ainda podem ser mesclados pelo BPE.
\end{exampleblock}

\vspace{4pt}
\centering
\textbf{Utilizado em:} GPT-2, GPT-3, GPT-4, Codex, LLaMA 2/3, \ldots

\end{frame}

% ============================================================

\begin{frame}{Byte-Level BPE vs.\ BPE Padrão}

\begin{center}
\begin{tikzpicture}[
    node distance=0.7cm,
    box/.style={draw, rounded corners, minimum width=3.2cm, minimum height=0.7cm, font=\small, align=center},
]
    % Standard BPE
    \node[font=\bfseries] (title1) {BPE Padrão};
    \node[box, fill=lightyellow, below=0.3cm of title1] (raw1) {Texto bruto (Unicode)};
    \node[box, fill=lightred, below=0.4cm of raw1] (base1) {Vocabulário base: chars Unicode\\(potencialmente 150K+)};
    \node[box, fill=lightblue, below=0.4cm of base1] (merge1) {BPE merges};
    \node[box, fill=lightgreen, below=0.4cm of merge1] (out1) {Pode gerar \texttt{<unk>}};
    
    \draw[->, thick] (raw1) -- (base1);
    \draw[->, thick] (base1) -- (merge1);
    \draw[->, thick] (merge1) -- (out1);
    
    % Byte-level BPE
    \node[font=\bfseries, right=3.5cm of title1] (title2) {Byte-Level BPE};
    \node[box, fill=lightyellow, below=0.3cm of title2] (raw2) {Texto bruto (Unicode)};
    \node[box, fill=lightgreen, below=0.4cm of raw2] (base2) {Vocabulário base: 256 bytes\\(fixo, completo)};
    \node[box, fill=lightblue, below=0.4cm of base2] (merge2) {BPE merges};
    \node[box, fill=lightgreen, below=0.4cm of merge2] (out2) {Sem \texttt{<unk>} possível};
    
    \draw[->, thick] (raw2) -- (base2);
    \draw[->, thick] (base2) -- (merge2);
    \draw[->, thick] (merge2) -- (out2);
\end{tikzpicture}
\end{center}

\vspace{4pt}

O GPT-2 usa um \textbf{mapeamento byte-para-unicode}: cada byte é representado como um caractere Unicode imprimível, de modo que o código BPE baseado em strings existente funciona sem modificações.

\end{frame}

% ============================================================
\section{Comparação e Considerações Práticas}
% ============================================================

\begin{frame}{OOV e Pré-tokenização}

\begin{columns}[T]
\column{0.48\textwidth}
\begin{block}{Tratamento de OOV (Out-Of-Vocabulary)}
\scriptsize
Tokens ou palavras que \textbf{não apareceram} no corpus de treinamento.
\begin{itemize}
    \item \textbf{Caractere:} cobertura total; o vocabulário são os próprios bytes/chars, nada é OOV
    \item \textbf{BPE / Unigram:} decompõe até chegar nos caracteres do vocabulário base (\textit{chars raros})
    \item \textbf{WordPiece:} substitui pelo token especial \texttt{[UNK]} se não conseguir segmentar
\end{itemize}
\end{block}

\column{0.48\textwidth}
\begin{block}{Pré-tokenização}
\scriptsize
Etapa \textbf{anterior} ao algoritmo: divide o texto em ``palavras'' por espaço e/ou pontuação.
\begin{itemize}
    \item \textbf{WordPiece (Sim):} o BERT sempre divide por espaço antes de aplicar o algoritmo
    \item \textbf{BPE (Geralmente):} o GPT-2 usa regras de pré-tokenização (ex.: nunca mescla letra com espaço)
    \item \textbf{SentencePiece (Opcional):} trata o espaço como caractere normal (\texttt{\_}) essencial para idiomas sem separação por espaço (japonês, chinês)
\end{itemize}
\end{block}
\end{columns}

\end{frame}


\begin{frame}{Comparação de Métodos de Tokenization}

\begin{center}
\small
\renewcommand{\arraystretch}{1.4}
\begin{tabular}{>{\bfseries}l c c c c}
\toprule
& \textbf{Caractere} & \textbf{BPE} & \textbf{WordPiece} & \textbf{Unigram} \\
\midrule
Direção            & --- & Bottom-up & Bottom-up & Top-down \\
Critério de merge  & --- & Frequência & Likelihood & Baseado em perda \\
Tratamento de OOV  & Total & Chars raros & \texttt{[UNK]} & Chars raros \\
Determinístico     & Sim & Sim & Sim & Não$^*$ \\
Pre-tokenization?  & Não & Geralmente & Sim & Opcional \\
\midrule
\multicolumn{5}{l}{\textbf{Modelos notáveis}} \\
\midrule
& ByT5 & Família GPT & BERT & T5, LLaMA \\
\bottomrule
\end{tabular}
\end{center}

\vspace{4pt}

{\footnotesize $^*$O Unigram pode produzir múltiplas segmentações válidas; durante o treinamento, isso pode servir como forma de \textbf{subword regularization} (data augmentation).}

\end{frame}

% ============================================================

\begin{frame}{Qual Tokenizer Cada Modelo Usa?}

\begin{center}
\small
\renewcommand{\arraystretch}{1.2}
\begin{tabular}{lll}
\toprule
\textbf{Modelo} & \textbf{Tokenizer} & \textbf{Tamanho do Vocabulário} \\
\midrule
BERT         & WordPiece                & 30,522 \\
GPT-2        & Byte-level BPE           & 50,257 \\
GPT-3        & Byte-level BPE (\texttt{r50k\_base})    & 50,257 \\
GPT-4        & Byte-level BPE (\texttt{cl100k\_base})  & 100,256 \\
T5           & SentencePiece (Unigram)  & 32,000 \\
LLaMA 1/2    & SentencePiece (BPE)       & 32,000 \\
LLaMA 3      & Byte-level BPE (tiktoken) & 128,256 \\
Gemma        & SentencePiece (BPE)      & 256,000 \\
Mistral      & SentencePiece (BPE)      & 32,000 \\
Claude       & Byte-level BPE (custom)  & --- \\
\bottomrule
\end{tabular}
\end{center}

\vspace{4pt}

\textbf{Tendência:} Modelos modernos estão migrando para vocabulários maiores e abordagens byte-level para melhorar a cobertura multilíngue.

\end{frame}

% ============================================================

\begin{frame}{Subword Regularization e BPE Dropout}

\textbf{Ideia:} A tokenization não precisa ser determinística durante o treinamento!

\vspace{6pt}

\begin{columns}[T]
\column{0.48\textwidth}
\begin{block}{Subword Regularization (Kudo, 2018)}
O Unigram pode amostrar entre múltiplas segmentações válidas:

\vspace{4pt}
\texttt{"lowest"} pode ser:
\begin{itemize}
    \item \tok{low}~\tok{est} (probabilidade 0.7)
    \item \tok{lo}~\tok{we}~\tok{st} (probabilidade 0.2)
    \item \tok{l}~\tok{o}~\tok{w}~\tok{est} (probabilidade 0.1)
\end{itemize}
\end{block}

\column{0.48\textwidth}
\begin{block}{BPE Dropout (Provilkov et al., 2020)}
Ignora aleatoriamente alguns merges do BPE durante o treinamento:

\vspace{4pt}
Com dropout $p=0.1$:
\begin{itemize}
    \item Normal: \tok{low}~\tok{est}
    \item Ignorado: \tok{lo}~\tok{w}~\tok{est}
    \item Ignorado: \tok{low}~\tok{e}~\tok{st}
\end{itemize}
\end{block}
\end{columns}

\vspace{6pt}

Ambos atuam como \textbf{data augmentation}, tornando os modelos mais robustos a segmentações raras.

\end{frame}

% ============================================================
\section{De Tokens a Embeddings}
% ============================================================

\begin{frame}[shrink=10]{De Token IDs a Embeddings}

O tokenizer produz uma sequência de \textbf{inteiros} (IDs). O transformer precisa de \textbf{vetores densos}.

\vspace{8pt}

\begin{columns}[T]
\begin{column}{0.48\textwidth}
\begin{block}{A tabela de embedding}
Uma matriz $E \in \mathbb{R}^{|V| \times d}$, onde cada linha é o vetor de um token:
\[
  \mathbf{e}_i = E[i] \in \mathbb{R}^d
\]
\texttt{nn.Embedding} é uma \textbf{tabela de lookup}, não uma transformação linear: apenas seleciona uma linha.
\end{block}
\end{column}
\begin{column}{0.48\textwidth}
\begin{center}
\renewcommand{\arraystretch}{1.2}
\small
\begin{tabular}{c|ccc}
\textbf{ID} & \multicolumn{3}{c}{$\mathbf{e} \in \mathbb{R}^d$} \\
\midrule
0   & $0.12$ & $-0.43$ & $\cdots$ \\
$\vdots$ & $\vdots$ & $\vdots$ & $\vdots$ \\
\rowcolor{lightblue}
\textbf{15496} & $\mathbf{0.31}$ & $\mathbf{-1.20}$ & $\mathbf{\cdots}$ \\
$\vdots$ & $\vdots$ & $\vdots$ & $\vdots$ \\
$|V|-1$ & $-0.07$ & $0.88$ & $\cdots$ \\
\end{tabular}

\vspace{4pt}
\small ID 15496 $\rightarrow$ linha 15496 $\rightarrow$ $\mathbf{e} \in \mathbb{R}^{768}$
\end{center}
\end{column}
\end{columns}

\vspace{6pt}

\begin{exampleblock}{Pipeline completo}
\centering
texto $\xrightarrow{\text{tokenizer}}$ [IDs] $\xrightarrow{\texttt{nn.Embedding}}$ [vetores] $\xrightarrow{+\,\text{PE}}$ entrada do transformer
\end{exampleblock}

\end{frame}

% ============================================================

\begin{frame}[fragile,shrink=10]{De Token IDs a Embeddings (código)}

\begin{exampleblock}{GPT-2 com HuggingFace + PyTorch}
\begin{verbatim}
from transformers import AutoTokenizer, AutoModelForCausalLM

tokenizer = AutoTokenizer.from_pretrained("gpt2")
model     = AutoModelForCausalLM.from_pretrained("gpt2")

ids  = tokenizer("Olá mundo", return_tensors="pt").input_ids
# ids.shape -> (1, seq_len)  -- tensor de inteiros

embs = model.transformer.wte(ids)   # nn.Embedding lookup
# embs.shape -> (1, seq_len, 768)  -- tensor de floats
\end{verbatim}
\end{exampleblock}

\vspace{6pt}

\begin{itemize}
    \item \texttt{wte} = \textit{word token embedding}, o nome usado internamente pelo GPT-2
    \item Após este passo, os embeddings posicionais são somados e a sequência entra na pilha de transformers
    \item O \texttt{model.transformer.wte} é uma instância de \texttt{torch.nn.Embedding(50257, 768)}
\end{itemize}

\end{frame}

% ============================================================

\begin{frame}[shrink=15]{Quantos Parâmetros Custam os Embeddings?}

\[
  \text{params}_{\text{embedding}} = |V| \times d
\]

\vspace{6pt}

\begin{center}
\small
\renewcommand{\arraystretch}{1.2}
\begin{tabular}{lrrrrr}
\toprule
\textbf{Modelo} & $|V|$ & $d$ & \textbf{Params emb.} & \textbf{Total} & \textbf{\% do total} \\
\midrule
GPT-2 Small  & 50,257 & 768   & 38.6M & 117M  & \textbf{33\%} \\
GPT-2 Medium & 50,257 & 1,024 & 51.5M & 345M  & \textbf{15\%} \\
GPT-2 Large  & 50,257 & 1,280 & 64.3M & 762M  & \textbf{8.4\%} \\
GPT-2 XL     & 50,257 & 1,600 & 80.4M & 1.5B & \textbf{5.4\%} \\
LLaMA 2 7B   & 32,000 & 4,096 & 131M    & 7B    & \textbf{1.9\%} \\
\bottomrule
\end{tabular}
\end{center}

\vspace{6pt}

\begin{alertblock}{Conclusões}
\begin{itemize}
    \item \textbf{Vocabulário maior = mais parâmetros}: escolhas de tokenization têm custo direto no tamanho do modelo
    \item \textbf{Modelos menores pagam mais}: em GPT-2 Small, $\frac{1}{3}$ dos parâmetros estão na tabela de embedding
    \item \textbf{Weight tying}: GPT-2 e muitos outros modelos \textbf{reutilizam} a mesma matriz como camada de saída (\texttt{lm\_head}), sem custo extra de parâmetros para a projeção final
\end{itemize}
\end{alertblock}

\end{frame}

% ============================================================
\section{Resumo}
% ============================================================

\begin{frame}{Principais Conclusões}

\begin{enumerate}
    \item \textbf{A tokenization é uma decisão arquitetural}, não uma etapa de pré-processamento

    \vspace{4pt}

    \item \textbf{Nível de caractere e de palavra} são extremos demais; métodos por subword encontram o ponto de equilíbrio

    \vspace{4pt}

    \item \textbf{BPE}: mescla os pares mais frequentes (bottom-up, baseado em frequência)

    \vspace{4pt}

    \item \textbf{WordPiece}: mescla pares maximizando a likelihood (bottom-up, baseado em probabilidade)

    \vspace{4pt}

    \item \textbf{Unigram}: começa grande, poda os tokens menos úteis (top-down, baseado em perda)

    \vspace{4pt}

    \item \textbf{SentencePiece}: um framework para tokenization agnóstica ao idioma (suporta BPE ou Unigram)

    \vspace{4pt}

    \item \textbf{Byte-level BPE}: garante cobertura total com 256 tokens base

\end{enumerate}

\end{frame}

% ============================================================
\section*{Referências e Leituras}
% ============================================================

\begin{frame}{Leituras Recomendadas}
    \leiturasize
    \begin{itemize}
        \item Radford, A. et al. (2019). Language Models are Unsupervised Multitask Learners. (GPT-2)
        \item Xue, L. et al. (2022). ByT5: Towards a Token-Free Future. \textit{NAACL}.
        \item Yu, L. et al. (2023). MegaByte: Predicting Million-Byte Sequences. \textit{NeurIPS}.
        \item Provilkov, I. et al. (2020). BPE-Dropout. \textit{ACL}.
    \end{itemize}
\end{frame}

\begin{frame}{Referências}

\small

\begin{itemize}
    \item Sennrich, R., Haddow, B., \& Birch, A. (2016). Neural Machine Translation of Rare Words with Subword Units. \textit{ACL}.
    \item Schuster, M. \& Nakajima, K. (2012). Japanese and Korean Voice Search. \textit{ICASSP}.
    \item Wu, Y. et al. (2016). Google's Neural Machine Translation System. \textit{arXiv:1609.08144}.
    \item Kudo, T. (2018). Subword Regularization. \textit{ACL}.
    \item Kudo, T. \& Richardson, J. (2018). SentencePiece. \textit{EMNLP}.
\end{itemize}

\end{frame}



\end{document}
